\begin{table}[htbp]\centering
{\fontsize{10}{12}\selectfont\renewcommand{\arraystretch}{1.18}\setlength{\tabcolsep}{4pt}
\begin{tabularx}{\linewidth}{@{}L{77mm}C{27mm}C{26mm}X@{}}\toprule\rowcolor{tablehead}
\textbf{Scenario} & \textbf{Purchase cost} & \textbf{Fee} & \textbf{Total outlay} \\ \midrule
One contract at \$0.50; direct precision & \$0.5000 & \$0.0175 & \$0.5175 \\ \arrayrulecolor{tablerule}\hline
One contract at \$0.50; cent precision & \$0.5000 & \$0.0200 & \$0.5200 \\ \arrayrulecolor{tablerule}\hline
Six one-contract legs at \$0.15; raw fee & \$0.90000 & \$0.05355 & \$0.95355 \\ \arrayrulecolor{tablerule}\hline
\arrayrulecolor{black}\bottomrule\end{tabularx}}
\caption[Illustrative fee arithmetic]{Illustrative fee arithmetic.}\label{tab:fees}
\vspace{2pt}{\centering\fontsize{9.3}{11.5}\selectfont Constructed examples with fee multiplier one. The last row is the raw-fee synthetic experiment, without fill-level balance rounding.\par}
\end{table}
